\chapter*{Förord}
\addcontentsline{toc}{chapter}{Förord}
\markboth{Förord}{Förord}

Ett hemsnickrat protokoll är lärorikt att bygga. Man bestämmer själv var framen börjar, hur lång
den är, hur den skyddas mot fel, och vad som ska hända när ett svar uteblir. Varje mekanism blir
ens egen --- och varje mekanism blir något som måste skrivas och underhållas.

Den här boken handlar om vad som händer när någon annan redan har bestämt allt det, en gång, och
gjutit det i kisel.

CAN är från 1986, det sitter i varje bil som byggts sedan mitten av nittiotalet, och det löser
nästan exakt de problem ni just löste själva --- fast i hårdvara, och bättre. Framing, checksumma,
kollisionshantering och kvittens ligger i kontrollern. Det som återstår för mjukvaran är att lämna
över en identifierare, en längd och några databytes, och sedan läsa av hur det gick.

\begin{aside}
\note Att CAN gör det åt er betyder inte att ni slipper förstå det. Bygger ni kontrollern är
protokollet er kravspecifikation, bit för bit. Skriver ni drivrutinen måste ni hantera precis det
kontrollern \emph{inte} kan säga er: varför en förlorad arbitrering ser ut som en lyckad
sändning, varför en frame som ingen kvitterar ändå fullbordas, och varför en mottagen frame kan
försvinna under näsan på er. Allt det är konsekvenser av hur CAN fungerar, och de står i den här
boken.
\end{aside}

\section*{Vad boken täcker}

\Kapref{ch:buss} börjar vid ledningarna: varför en delad buss, vad dominanta och recessiva bitar
är, och varför den asymmetrin är hela grunden. \Kapref{ch:frameformat} går igenom CAN-framen fält för fält.
\Kapref{ch:stuffing} tar bitstoppningen och CRC-15:an, de två mekanismerna som gör att en frame
överlever en störig buss. \Kapref{ch:arbitrering} arbetar igenom arbitreringen bit för bit, för
hand. \Kapref{ch:tajming} handlar om bittajming och sampelpunkten: hur noder utan gemensam klocka
ändå läser samma bit samtidigt. \Kapref{ch:fel} beskriver felhanteringen i riktig CAN --- felframes,
felräknare och bus-off.

\Kapref{ch:kursen} är bokens enda kapitel som inte handlar om CAN i allmänhet. Det beskriver den
förenklade kontroller som byggs i undervisningen, vad den implementerar av allt det föregående,
och vad den medvetet låter bli. Det kapitlet är det som kommer till användning vid felsökning ---
från båda hållen.

\section*{Vad boken inte täcker}

Boken beskriver \textbf{klassisk CAN} med 11-bitars identifierare, alltså CAN 2.0A. Utökade
29-bitars identifierare (2.0B) nämns där de påverkar framens form, men arbetas inte igenom. CAN~FD
och CAN~XL ligger utanför, liksom högre lager som CANopen och J1939.

Boken beskriver inte heller elektriken i detalj: differentiell signalering, termineringens
fysik och kabellängdens inverkan berörs bara så långt det behövs för att förstå varför bussen
ser ut som den gör.

\section*{Övningar}

Varje kapitel slutar med några övningar. De görs med papper och penna --- det finns ingen kod att
kompilera i den här boken --- och svaren på samtliga står i Appendix~A.

Övningarna är korta med avsikt. De är till för att kontrollera att kapitlet fastnade, inte för
att ta en kväll.
